% !TEX program = xelatex
\documentclass[12pt,a4paper]{ctexart}

% ==================== 页边距设置 ====================
\usepackage[top=2.5cm, bottom=2.5cm, left=3cm, right=2.5cm]{geometry}

% ==================== 全局间距压缩 ====================
\usepackage{enumitem}
\setlist{nosep, leftmargin=*, itemsep=0pt, parsep=0pt, topsep=0pt, partopsep=0pt}
\setlength{\parskip}{0pt}
\setlength{\parindent}{2em}
\setlength{\abovedisplayskip}{3pt plus 1pt minus 1pt}
\setlength{\belowdisplayskip}{3pt plus 1pt minus 1pt}
\setlength{\abovedisplayshortskip}{0pt plus 1pt}
\setlength{\belowdisplayshortskip}{2pt plus 1pt minus 1pt}
\setlength{\intextsep}{3pt plus 1pt minus 1pt}
\setlength{\textfloatsep}{6pt plus 2pt minus 2pt}
\setlength{\floatsep}{6pt plus 2pt minus 2pt}
\setlength{\dbltextfloatsep}{6pt plus 2pt minus 2pt}
\setlength{\dblfloatsep}{6pt plus 2pt minus 2pt}
\setlength{\abovecaptionskip}{3pt plus 1pt minus 1pt}
\setlength{\belowcaptionskip}{0pt}
\AtBeginEnvironment{table}{\setlength{\tabcolsep}{4pt}\renewcommand{\arraystretch}{1.05}}
\AtBeginEnvironment{tabular}{\small}
\ctexset{
    section = {
        beforeskip = {6pt plus 2pt minus 2pt},
        afterskip = {4pt plus 2pt minus 2pt},
    },
    subsection = {
        beforeskip = {4pt plus 2pt minus 2pt},
        afterskip = {2pt plus 1pt minus 1pt},
    },
    subsubsection = {
        beforeskip = {3pt plus 1pt minus 1pt},
        afterskip = {1pt plus 1pt minus 1pt},
    },
}

% ==================== 字体与编码 ====================
\usepackage{fontspec}
\setmainfont{Times New Roman}
\setmonofont{Courier New}

% ==================== 中文支持 ====================
\usepackage{xeCJK}
\setCJKmainfont{SimSun}[AutoFakeBold=true]

% ==================== 图表与图形 ====================
\usepackage{graphicx}
\usepackage{float}
\usepackage{booktabs}
\usepackage{longtable}
\usepackage{array}
\usepackage{colortbl}
\usepackage{xcolor}

% ==================== 数学与算法 ====================
\usepackage{amsmath, amssymb, amsthm}
\usepackage[ruled,linesnumbered]{algorithm2e}

% ==================== 超链接与引用 ====================
\usepackage[colorlinks=true,linkcolor=blue,citecolor=blue,urlcolor=blue]{hyperref}

% ==================== 文献引用 ====================
\usepackage{cite}

% ==================== 代码高亮 ====================
\usepackage{listings}
\lstset{
    basicstyle=\small\ttfamily,
    breaklines=true,
    frame=single,
    numbers=left,
    numberstyle=\tiny,
    xleftmargin=2em,
    framexleftmargin=1.5em,
    language=Python,
    showstringspaces=false,
    keywordstyle=\color{blue},
    commentstyle=\color{gray},
    stringstyle=\color{red}
}

% ==================== 页面样式 ====================
\usepackage{fancyhdr}
\pagestyle{fancy}
\fancyhf{}
\fancyhead[L]{\small 数据挖掘课程设计开题报告}
\fancyhead[R]{\small 基于建筑能耗对标数据的节能改造优先级评估}
\fancyfoot[C]{\thepage}
\setlength{\headheight}{14pt}
\renewcommand{\headrulewidth}{0.4pt}
\renewcommand{\footrulewidth}{0pt}

% ==================== 标题格式 ====================
\ctexset{
    section = {
        format = {\Large\bfseries\centering},
        name = {,、},
        number = \chinese{section},
    },
    subsection = {
        format = {\large\bfseries},
        name = {,、},
        number = \arabic{subsection},
    }
}

% ==================== 文档信息 ====================
\title{{\Huge\bfseries 数据挖掘\\[10pt] 课程设计开题报告}}
\author{}
\date{}

\raggedbottom

\begin{document}

% ==================== 封面 ====================
\maketitle

\vspace{1.5cm}

\begin{center}
    {\LARGE\bfseries 基于建筑能耗对标数据的\\[8pt] 城市建筑能耗预测与节能改造优先级评估}

    \vspace{2cm}

    {\large\begin{tabular}{rl}
        \textbf{课题名称：} & 基于西雅图建筑能耗对标数据的城市建筑能耗预测与节能改造优先级评估 \\
        \textbf{课程名称：} & 数据挖掘与大数据分析 \\
        \textbf{完成日期：} & 2026年7月 \\
    \end{tabular}}
\end{center}

\vspace{1cm}
\begin{center}
    \rule{0.6\textwidth}{0.5pt}
\end{center}

\newpage

% ==================== 目录 ====================
\tableofcontents
\newpage

% ==================== 正文 ====================

\section{研究背景与问题定义}

\subsection{研究背景}

建筑行业能耗占全球终端能源消费约30\%，碳排放占总量约28\%\cite{unep2020global}。中国在"双碳"目标（2030年碳达峰、2060年碳中和）背景下，既有建筑存量巨大且能耗水平参差不齐，如何精准识别高耗能建筑并制定科学的节能改造初步筛选清单，是城市管理者面临的核心挑战。

本研究依托西雅图市建筑能耗对标数据集（Seattle Building Energy Benchmarking, 2015--2024），该数据以``建筑--年份''为观测单元共计34,699条记录，含47个字段，覆盖能耗强度（SiteEUI）、温室气体排放（GHG）、建筑物理属性（建筑面积、层数、建造年份）、能源之星评分等信息。同时引入西雅图市按建筑类型和年份汇集的基准性能范围数据（Benchmarking Performance Ranges by Building Type），为能耗对标评估提供行业参照。

\subsection{问题定义}

城市建筑节能改造初步筛选面临三个核心问题：
\begin{enumerate}
    \item \textbf{能耗基准估计问题：}仅基于建筑静态属性（类型、面积、建造年代、区域等）和年度背景信息等预测时可获得的数据，能否合理估计建筑在特定年度的SiteEUI基准水平？
    \item \textbf{异常高耗能建筑识别问题：}哪些建筑实际能耗显著高于同类建筑的合理预期（模型预测基准和同类型统计基准）？这些建筑是否在近三年持续高耗能？
    \item \textbf{优先级排序问题：}在缺乏详细改造成本数据的条件下，如何综合预测残差、同类基准超耗、持续高耗能表现、碳排放强度、建筑老化程度等多维因素，构建建筑级节能改造优先级评分，并通过稳定性分析验证排序的可靠性？
\end{enumerate}

\subsection{数据说明}

本研究所用数据为公开的西雅图建筑能耗对标数据集，以"建筑--年份"为观测单元，共34,699条建筑年度记录，唯一建筑数量需逐年统计。同时使用西雅图市发布的按建筑类型与年份划分的基准性能范围数据，为计算相对基准的超耗比例提供参照。



\section{数据来源、字段说明与文献依据}

\subsection{数据来源}

本研究使用两个公开数据集：

1. \textbf{主数据集：}西雅图建筑能耗对标数据（Building Energy Benchmarking Data, 2015--Present），由西雅图市政府规划与社区发展办公室公开发布，包含34,699条"建筑--年份"观测记录，47个字段\cite{seattle_benchmarking}。

2. \textbf{基准数据集：}按建筑类型与年份划分的基准性能范围数据（Benchmarking Performance Ranges by Building Type），提供各类建筑能耗的百分位参照\cite{seattle_benchmarking}。

原始数据总规模约15 MB，适合在个人计算机上处理。

\subsection{关键字段分类}

主数据集含47个字段，涵盖建筑标识（ID、名称、地址）、静态属性（类型、面积、建造年份）、能源消耗（场地能耗、源能耗、EUI）、温室气体排放（总量与强度）以及ENERGY STAR评分。为避免目标泄漏并明确建模边界，所有字段按以下六类划分，分别制定使用策略：

\begin{table}[H]
\centering
\caption{字段分类与使用策略}
\label{tab:field_classification}
\begin{tabular}{p{2.2cm}p{3.8cm}p{5.5cm}}
\hline
\textbf{类别} & \textbf{示例字段} & \textbf{使用策略} \\
\hline
建筑静态属性 & 建筑类型、建筑面积(GFA)、层数、建造年份、经纬度 & 保留为模型特征 \\
\hline
时间特征 & DataYear、BuildingAge & 保留，用于年份切分及建筑年龄计算 \\
\hline
建筑用途 & LargestPropertyUseType、Second/ThirdLargestPropertyUseType & 编码后作为特征 \\
\hline
能耗目标 & SiteEUI (kBtu/sf) & 回归目标变量 \\
\hline
同期能耗派生 & SourceEUI、SiteEnergyUse、SourceEnergyUse、SiteEUIWN、SourceEUIWN & 剔除——与目标使用相同年度能耗信息，预测时不可提前获取 \\
\hline
同期排放变量 & TotalGHGEmissions、GHGEmissionsIntensity & 预测模型中剔除，排序阶段可作为参考 \\
\hline
同期评分 & ENERGYSTARScore & 单独评估——作为缺失指标（ES\_Missing）参与建模，不直接作为特征 \\
\hline
标识字段 & OSEBuildingID、建筑名称、地址 & 不进入模型，仅用于标识 \\
\hline
\end{tabular}
\end{table}

\subsection{相关研究方法}

\begin{table}[H]
\centering
\caption{相关研究对比}
\label{tab:literature}
\begin{tabular}{p{1.8cm}p{2.5cm}p{2.8cm}p{2.5cm}p{2.5cm}}
\hline
\textbf{方法类别} & \textbf{代表方法} & \textbf{数据类型} & \textbf{主要指标} & \textbf{局限} \\
\hline
工程法(白箱) & EnergyPlus, eQUEST & 建筑物理参数 & 模拟误差 & 参数要求高，难以大规模应用 \\
\hline
统计法(灰箱) & 线性回归、时间序列 & 历史能耗数据 & RMSE、MAE \cite{ashrae_energy} & 非线性关系捕捉能力有限 \\
\hline
集成学习(黑箱) & 随机森林、XGBoost、LightGBM \cite{rf_breiman}\cite{xgboost_chen}\cite{lgbm_ke} & 建筑对标数据 & $R^2$、MAE & 可解释性不足 \\
\hline
模型解释 & SHAP \cite{shap_lundberg} & 模型输出 & 特征重要性 & 不能解释为因果关系 \\
\hline
综合评价 & AHP、TOPSIS \cite{ahp_saaty} & 多指标数据 & 综合得分 & 权重主观性较强 \\
\hline
\end{tabular}
\end{table}

现有研究多聚焦于单一环节（仅预测或仅评价），将能耗预测、模型解释、建筑分群和改造候选排序整合为完整分析链的工作仍较少，这也是本课题的主要工作方向。

\subsection{基准数据的应用}

本研究将基准性能范围数据纳入分析流程，用于：
\begin{enumerate}
    \item 计算每栋建筑相对于同类型同年度基准SiteEUI的超耗比例：\[ \text{超耗比例} = \frac{\text{实际 SiteEUI} - \text{同类型年度基准 SiteEUI}}{\text{同类型年度基准 SiteEUI}} \]
    \item 为优先级排序提供更客观的参照，避免医院、酒店等天然高能耗建筑因绝对SiteEUI值较高而全部位列前茅。
\end{enumerate}

\section{研究目标与评价指标}

\subsection{数据挖掘问题定义}

本项目围绕西雅图市建筑能耗对标数据，解决以下数据挖掘问题：

\begin{table}[H]
\centering
\caption{数据挖掘问题类型与任务映射}
\label{tab:dm_problems}
\begin{tabular}{p{2.5cm}p{3cm}p{6.5cm}}
\hline
\textbf{问题类型} & \textbf{对应阶段} & \textbf{说明} \\
\hline
数据清洗与预处理 & 阶段一 & 跨年度字段审计、缺失值填补、异常值检测与处理、目标泄漏字段识别与剔除 \\
\hline
描述性统计与可视化 & 阶段二 & SiteEUI分布分析、建筑类型能耗对比、特征相关性分析、缺失模式探索 \\
\hline
监督学习—回归 & 阶段三 & 以建筑静态属性为特征，估计建筑在特定年度的合理SiteEUI基准水平；比较多种回归模型，通过滚动时间窗口交叉验证评估泛化性能 \\
\hline
模型可解释性分析 & 阶段四 & 利用SHAP方法识别对模型预测有重要贡献的特征，分析特征与预测值之间的关联模式 \\
\hline
异常检测 & 阶段五 & 基于标准化预测残差和同类统计基准，三层递进筛选异常高耗能建筑；通过持续性验证减少偶然性误判 \\
\hline
多指标综合排序 & 阶段六 & 构建多维加权评分体系，对建筑进行节能改造优先级排序，并通过多种方案对比和重采样验证排序稳定性 \\
\hline
聚类分析（辅助） & 阶段七 & 对高优先级候选建筑进行聚类画像，辅助理解候选建筑的类型结构和特征模式 \\
\hline
\end{tabular}
\end{table}

\subsection{总体目标}

以西雅图市2015--2024年建筑能耗对标数据为基础，利用建筑类型、用途、面积、层数、建造年份、区域和年度背景等预测时可获得的信息，估计建筑在特定年度的合理SiteEUI基准水平；在此基础上识别异常高耗能建筑，汇总建筑近三年持续表现，构建``数据审计 $\rightarrow$ 能耗基准预测 $\rightarrow$ 异常高耗能识别 $\rightarrow$ 建筑级汇总 $\rightarrow$ 改造优先级排序 $\rightarrow$ 辅助聚类画像''的完整分析链，为城市既有建筑节能改造初步筛选提供数据驱动的决策参考。

\subsection{各任务输入、输出与评价指标}

\begin{table}[H]
\centering
\caption{研究任务说明}
\label{tab:task_io}
\begin{tabular}{p{0.8cm}p{2.5cm}p{2.5cm}p{2.5cm}p{3.5cm}}
\hline
\textbf{序号} & \textbf{任务名称} & \textbf{输入} & \textbf{输出} & \textbf{评价指标} \\
\hline
T1 & 能耗基准回归预测 & 建筑静态属性特征矩阵、SiteEUI目标向量 & 最优回归模型及预测值 & MAE、RMSE、$R^2$；目标为显著优于全体均值基线和建筑类型中位数基线（如MAE降低10\%--20\%），不设刚性$R^2$阈值 \\
\hline
T2 & 重要预测特征识别 & 最佳模型、测试集 & SHAP特征重要性排名与可视化图表 & SHAP全局重要性排序、依赖图、瀑布图；高度相关特征分组解释；统一使用``重要预测特征''表述 \\
\hline
T3 & 异常高耗能建筑识别 & 模型预测值、实际SiteEUI、同类基准数据 & 异常建筑候选列表 & 三层递进筛选（标准化残差$>1.5$ + 同类P75基准超标 + 近三年持续出现）；报告各层筛选通过数量与重叠率 \\
\hline
T4 & 建筑级记录汇总与优先级评分 & 年度建筑记录、异常识别结果 & 建筑级优先级排名表 & 6维度加权评分；数据质量折扣修正；Top-N清单在不同权重方案下的重合率与Jaccard相似度 \\
\hline
T5 & 排序稳定性检验 & 多权重方案下的排序结果 & 稳定性分析报告 & Spearman/Kendall秩相关系数、Bootstrap入选频率、各建筑类型候选占比合理性 \\
\hline
T6 & 候选建筑聚类画像 & 高优先级候选建筑特征 & 聚类画像报告与可视化 & 轮廓系数、Calinski-Harabasz指数、画像可解释性 \\
\hline
\end{tabular}
\end{table}

\subsection{评估基线}

为验证机器学习模型的附加值，设置以下基线模型：

\begin{enumerate}
    \item \textbf{基线A（全体均值）：}以训练集SiteEUI均值作为所有建筑的预测值；
    \item \textbf{基线B（建筑类型中位数）：}按建筑类型分组，以训练集各组SiteEUI中位数作为同类建筑的预测值。
\end{enumerate}

只有当复杂模型在测试集上显著优于这些简单基线时，模型才具有实用价值。

\subsection{预测任务定义}

本项目将预测任务统一定义为：

\begin{quote}
利用建筑类型、用途、面积、层数、建造年份、区域和年度背景等\textbf{预测时可获得的信息}，估计建筑在特定年度的\textbf{合理SiteEUI基准水平}。
\end{quote}

该定义明确了项目的性质：监督学习回归与建筑能耗对标（benchmarking），而非严格的时间序列预测。除非后续明确加入上一年SiteEUI、近三年均值、同比变化率等滞后特征，否则不使用``历史能耗预测下一年能耗''的表述。


\section{数据预处理与泄漏控制}

\subsection{跨年度数据审计}

正式分析前，先对2015--2024年共10年的原始数据进行系统性审计：

\begin{table}[H]
\centering
\caption{跨年度数据审计内容}
\label{tab:audit}
\begin{tabular}{p{3.5cm}p{8.5cm}}
\hline
\textbf{审计项目} & \textbf{检查内容与输出} \\
\hline
字段名称一致性 & 各年份CSV的列名是否变化（如大小写、空格、括号差异），生成字段映射表 \\
\hline
字段单位一致性 & 同一字段在不同年份的单位是否一致（如kBtu vs GJ），生成字段映射表中的单位映射 \\
\hline
建筑类型命名一致性 & BuildingType和LargestPropertyUseType的取值在不同年份是否统一，是否存在拼写差异或分类合并 \\
\hline
建筑ID稳定性 & OSEBuildingID是否跨年稳定，是否存在同一ID在不同年份指向不同建筑的情况 \\
\hline
重复记录检查 & 同一建筑在同一年份是否有多条记录，完全重复行比例 \\
\hline
数值字段解析 & 是否存在因格式问题（如千分位逗号、特殊字符）导致数值字段被解析为字符串的情况 \\
\hline
各字段覆盖率 & 每年各字段的非缺失比例，生成字段覆盖率统计表 \\
\hline
PropertyGFATotal核查 & 重新检查该字段在各年份的实际覆盖率，确认是否因字段名称变化导致误判为全部缺失 \\
\hline
\end{tabular}
\end{table}

\subsection{数据预处理流程}

数据预处理按以下顺序执行，所有拟合操作（填充值、编码映射、标准化参数）仅在训练集上计算：

\begin{table}[H]
\centering
\caption{预处理步骤说明}
\label{tab:preprocessing}
\begin{tabular}{p{1cm}p{3.5cm}p{7.5cm}}
\hline
\textbf{步骤} & \textbf{操作} & \textbf{说明} \\
\hline
P1 & 字段统一 & 去除列名中括号、空格、斜杠等特殊字符；统一大小写；修正各年份字段名差异 \\
\hline
P2 & 数值类型转换 & 将应为数值型的列强制转换为numeric，无法解析的值标记为NaN \\
\hline
P3 & 唯一建筑与重复检查 & 同年同OSEBuildingID重复、完全重复行检查；记录重复条数与处理方式 \\
\hline
P4 & 目标泄漏审查 & 识别并剔除与SiteEUI存在确定性数量关系或共享年度能耗测量信息的派生字段（SourceEUI、SiteEnergyUse、SourceEnergyUse、天气标准化字段等） \\
\hline
P5 & 目标变量过滤 & 仅删除SiteEUI缺失、SiteEUI $\le$ 0和明确录入错误的记录；\textbf{不进行分位数截尾}，保留合理的极端高能耗建筑作为后续分析对象 \\
\hline
P6 & YearBuilt处理 & YearBuilt $>$ DataYear的异常值标记为缺失，增加异常指示变量YearBuilt\_Anomaly；按建筑类型中位年份填补；确保BuildingAge非负 \\
\hline
P7 & 缺失值填充 & 数值字段优先以同建筑类型中位数填充，仍缺失时使用训练集整体中位数；类别字段填充为``Unknown''；同时生成各字段的缺失指示变量（如ENERGYSTARScore\_missing） \\
\hline
P8 & 特征工程 & 计算BuildingAge、LogGFA、面积比率、混合用途标识、缺失指示变量等（详见第5节特征列表） \\
\hline
P9 & 类别编码 & ElasticNet使用OneHotEncoder；树模型使用OneHotEncoder（LightGBM可选用原生类别特征）；\textbf{禁止对建筑类型使用LabelEncoder}以避免虚假的大小关系 \\
\hline
P10 & 目标变换与标准化 & 建模时对目标取$y = \log(1+\text{SiteEUI})$，预测后还原$\hat{y}_{\text{orig}} = \exp(\hat{y}_{\text{log}})-1$；RobustScaler在训练集上拟合，对验证集和测试集仅transform \\
\hline
\end{tabular}
\end{table}

\subsection{目标泄漏审查}

目标泄漏是建筑能耗建模中最关键的风险控制措施。在预测SiteEUI时，任何从同期能耗信息中派生的字段若进入特征集，都将导致模型在训练时``看到''目标的部分信息，从而高估真实预测能力。本课题在建模前对所有字段逐一审查，确保以下三类问题不会发生：

\textbf{第一类：直接泄漏——由目标变量数学变换得到的字段。}SourceEUI由SiteEUI乘以源-站转换系数计算得出（$\text{SourceEUI} = \text{SiteEUI} \times \text{Source-to-Site Factor}$）；SiteEnergyUse为建筑面积 $\times$ SiteEUI。这些字段与目标存在确定性的数量关系，一经发现立即剔除。

\textbf{第二类：间接泄漏——与目标共享年度能耗测量信息的字段。}SourceEnergyUse、SiteEnergyUseWN、SourceEUIWN等字段虽然计算方式不同，但均来自同一份年度能耗账单数据。在预测场景下（使用静态属性估计合理能耗水平），这些字段的值无法提前获知，因此必须剔除。

\textbf{第三类：灰色变量——与目标相关但非直接派生的字段。}TotalGHGEmissions和GHGEmissionsIntensity由能耗乘以排放因子计算，预测模型中不可使用，但在后续的优先级排序阶段可作为减碳维度的参考指标。ENERGYSTARScore实际由EPA基于当年SiteEUI百分位计算，具有较强的目标相关性，因此主模型（模型A）不使用，仅编码为二值指示变量（ES\_Missing）；模型B使用ENERGYSTARScore仅作为扩展对比实验。

\begin{table}[H]
\centering
\caption{目标泄漏审查与字段处理决策}
\label{tab:leakage}
\begin{tabular}{p{3.5cm}p{2.5cm}p{6cm}}
\hline
\textbf{字段} & \textbf{处理} & \textbf{理由} \\
\hline
SourceEUI、SourceEUIWN & 剔除 & 由SiteEUI乘以源-站转换系数计算，与目标存在确定性数量关系 \\
\hline
SiteEnergyUse、SiteEnergyUseWN、SourceEnergyUse & 剔除 & 能耗总量 $\div$ 面积可直接还原为SiteEUI，属直接泄漏 \\
\hline
TotalGHGEmissions、GHGEmissionsIntensity & 预测模型剔除，排序可用 & 由能耗和排放因子计算，预测时不可预先获得；排序阶段作为减碳参考 \\
\hline
ENERGYSTARScore & 主模型不直接入模，仅保留ES\_Missing指示变量；模型B使用作扩展对比 & 由EPA基于当年SiteEUI百分位计算，与目标高度相关；防止模型因评级信息而虚高 \\
\hline
PropertyGFATotal & 需重新核查各年份实际覆盖率 & 不应直接假定全部缺失，可能因字段名称变化导致匹配失败 \\
\hline
\end{tabular}
\end{table}

\subsection{SiteEUI异常值处理策略}

高能耗建筑正是本研究的识别对象，因此\textbf{不能对SiteEUI简单分位数截尾}。处理策略为：

\begin{enumerate}
    \item 删除SiteEUI缺失、非正值和明确录入错误（如负数、数量级异常的错误值）的记录；
    \item \textbf{保留合理的极端高能耗记录}，它们将在后续的异常识别和优先级排序中被标记和分析；
    \item 建模时对目标取对数$y = \log(1+\text{SiteEUI})$以缓解右偏分布对模型训练的影响；
    \item 预测后通过$\hat{y}_{\text{orig}} = \exp(\hat{y}_{\text{log}}) - 1$恢复为原始尺度；
    \item 原始SiteEUI继续用于异常识别和优先级排序。
\end{enumerate}

\subsection{建筑类型—年份基准匹配}

利用基准性能范围数据（Benchmarking Performance Ranges by Building Type）：

\begin{enumerate}
    \item 为每栋建筑匹配同类型、同年度的SiteEUI基准值（如P50中位数和P75上四分位数）；
    \item 计算超耗比例：正值表示高于同类平均水平，负值表示低于平均水平；
    \item 在优先级评分阶段，使用同类型超耗比例替代绝对SiteEUI，降低建筑类型固有差异的干扰，避免办公楼、医院等类型因天然能耗高而垄断Top-50。
\end{enumerate}

\subsection{目标变换}

SiteEUI通常呈右偏分布。本课题在训练时对目标进行$\log(1+\text{SiteEUI})$变换以改善残差分布的对称性，评估时在原始尺度（kBtu/sf）上报告误差。

\subsection{划分策略}

采用按年份的扩展窗口滚动验证（详见第5节），不再使用TimeSeriesSplit的机械划分。同一建筑可能出现在多个集合中（因为同栋建筑在不同年份均有记录），这符合``利用历史数据估计基准''的场景设定。辅助实验采用按OSEBuildingID的GroupShuffleSplit，检验模型对完全未见建筑的泛化能力。

\subsection{技术可行性}

本研究的技术基础坚实：

\begin{enumerate}
    \item \textbf{数据可用性：}西雅图市政府公开发布，数据完整规范，附带元数据说明，覆盖2015--2024共10年；
    \item \textbf{工具链成熟度：}Python Scikit-learn、XGBoost、LightGBM、SHAP均为工业级稳定库；
    \item \textbf{计算资源：}数据规模约15 MB，在合理控制参数搜索范围的情况下，可在课程项目周期内完成全部训练与解释分析；
    \item \textbf{方法论支撑：}回归预测、SHAP可解释性、K-Means聚类、多维加权评分均为经过充分验证的方法，本研究在此基础上进行整合与改进；
    \item \textbf{可复现性：}全流程Python脚本化，环境依赖写入requirements.txt；
    \item \textbf{经验基础：}已完成鸢尾花数据集分析项目，具备完整的ML pipeline实施经验。
\end{enumerate}

\subsection{主要难点与应对}

\begin{table}[H]
\centering
\caption{技术难点与应对策略}
\label{tab:challenges}
\begin{tabular}{p{3cm}p{4cm}p{5cm}}
\hline
\textbf{难点} & \textbf{原因} & \textbf{应对策略} \\
\hline
缺失值处理 & 部分字段缺失率达20--40\% & 数值以同建筑类型中位数优先填充，类别填充``Unknown''，同时生成缺失指示变量保留缺失信息 \\
\hline
不同建筑类型样本量差异大 & 办公楼远多于医院等类型 & 评估时按建筑类型分组报告性能，排序使用同类型超耗比例避免类型垄断 \\
\hline
目标泄漏风险 & 派生变量与目标共享能耗信息 & 严格剔除同期派生字段，仅保留预测时可获得的静态属性 \\
\hline
信息泄漏于测试集 & 缺失值填充等操作在全量数据执行 & 所有拟合（中位数、编码器、标准化器）仅在训练集上计算 \\
\hline
高能耗建筑极端值 & 无法区分录入错误与真实高耗能 & 采用log1p变换缓解偏度；结合预测残差和同类基准双重验证 \\
\hline
K值选择主观 & 不同K值给出不同分群结果 & 综合轮廓系数+Calinski-Harabasz指数+画像可解释性综合判断 \\
\hline
权重系数主观 & 优先级评分中各维度重要性缺乏客观标准 & 设计多个权重方案（节能优先/减碳优先/大型建筑优先），通过Bootstrap、Jaccard相似度、Kendall相关系数检验排序稳定性 \\
\hline
PropertyGFATotal高缺失 & 可能因字段名称变化导致匹配失败 & 跨年度审计时重新核查，若确认为真实缺失则使用PropertyGFABuildings替代 \\
\hline
\end{tabular}
\end{table}


\section{模型设计与评估方案}

\subsection{总体技术流程}

本项目以一条清晰主线组织完整分析流程：

\begin{center}
\begin{tabular}{c}
\hline
\textbf{跨年度字段和数据质量审计} \\
$\downarrow$ \\
\textbf{按年份划分训练、验证和测试集} \\
$\downarrow$ \\
\textbf{构造无数据泄漏的建筑属性特征} \\
$\downarrow$ \\
\textbf{建立建筑能耗基准回归模型（多模型对比）} \\
$\downarrow$ \\
\textbf{时间外测试（2023、2024）和未见建筑测试（GroupShuffleSplit）} \\
$\downarrow$ \\
\textbf{SHAP解释重要预测特征} \\
$\downarrow$ \\
\textbf{基于标准化残差和同类基准识别异常高耗能建筑} \\
$\downarrow$ \\
\textbf{按OSEBuildingID汇总最近三年表现为建筑级记录} \\
$\downarrow$ \\
\textbf{构建建筑级改造优先级评分并检验排序稳定性} \\
$\downarrow$ \\
\textbf{对高优先级候选建筑进行辅助聚类画像} \\
\hline
\end{tabular}
\end{center}

\subsection{特征工程}

\subsubsection{推荐输入特征}

从原始字段中构造以下特征用于能耗基准预测：

\begin{table}[H]
\centering
\caption{模型特征列表}
\label{tab:features}
\begin{tabular}{p{5cm}p{7cm}}
\hline
\textbf{特征} & \textbf{计算方式与说明} \\
\hline
BuildingType & 建筑类型，使用OneHotEncoder编码 \\
\hline
LargestPropertyUseType & 主要用途类型，使用OneHotEncoder编码 \\
\hline
PropertyGFATotal & 总建筑面积（需重新核查该字段在各年份的实际覆盖率，不应直接判定为全部缺失） \\
\hline
NumberofFloors & 楼层数 \\
\hline
NumberofBuildings & 建筑数量 \\
\hline
YearBuilt & 建造年份（异常值标记后按建筑类型中位数填补） \\
\hline
BuildingAge & $\max(\text{DataYear} - \text{YearBuilt}, 0)$ \\
\hline
Neighborhood & 所在区域 \\
\hline
ParkingRatio & $\text{PropertyGFAParking} / \text{GFA}$ \\
\hline
AreaPerFloor & $\text{GFA} / \text{NumberofFloors}$ \\
\hline
AreaPerBuilding & $\text{GFA} / \text{NumberofBuildings}$ \\
\hline
LogGFA & $\log(1 + \text{GFA})$ \\
\hline
IsMixedUse & 建筑用途类型数 $> 1$ 的指示变量 \\
\hline
PrimaryUseAreaRatio & 主要用途面积占GFA的比例 \\
\hline
DataYear & 数据记录年份（作为年度背景变量） \\
\hline
ES\_Missing & ENERGYSTARScore是否缺失的二值指示变量 \\
\hline
YearBuilt\_Anomaly & YearBuilt异常的指示变量（YearBuilt $>$ DataYear 标记） \\
\hline
\end{tabular}
\end{table}

\subsubsection{谨慎使用或剔除的字段}

当目标为SiteEUI时，主模型不得使用以下字段：
\begin{itemize}
    \item 当年SiteEnergyUse、SourceEUI、天气标准化EUI（由目标直接派生或换算）；
    \item 当年电力、天然气、蒸汽消耗量（与目标共享年度能耗测量信息）；
    \item 当年TotalGHGEmissions、GHGEmissionsIntensity（由能耗乘以排放因子计算，预测时不可预先获得）；
    \item 任何与目标存在确定性数量关系的字段。
\end{itemize}

GHG排放强度仅在优先级排序阶段作为减碳维度参考。

\subsubsection{ENERGY STAR Score处理}

设置两组实验：
\begin{itemize}
    \item \textbf{模型A（主结果）：}不使用ENERGY STAR Score，仅保留ES\_Missing指示变量；
    \item \textbf{模型B（扩展对比）：}使用ENERGY STAR Score作为特征，仅用于对比模型A，检验评分信息对预测性能的提升幅度。
\end{itemize}

\subsubsection{类别编码策略}

不同模型采用不同的编码方案，禁止对建筑类型使用LabelEncoder：
\begin{itemize}
    \item \textbf{ElasticNet / 线性回归：}OneHotEncoder + RobustScaler；
    \item \textbf{随机森林 / XGBoost：}缺失填补 + OneHotEncoder；
    \item \textbf{LightGBM：}可OneHotEncoder，也可使用原生类别特征（dtype="category"）；
    \item \textbf{K-Means聚类：}仅使用连续数值特征，建筑类型不参与距离计算，仅用于聚类结果的事后解释。
\end{itemize}

\subsection{训练、验证与测试划分}

\subsubsection{主实验：按完整年份扩展窗口}

\begin{table}[H]
\centering
\caption{按年份滚动验证方案}
\label{tab:split}
\begin{tabular}{p{3cm}p{3cm}p{2.5cm}}
\hline
\textbf{训练集} & \textbf{验证集} & \textbf{用途} \\
\hline
2015--2017 & 2018 & 交叉验证折1 \\
2015--2018 & 2019 & 交叉验证折2 \\
2015--2019 & 2020 & 交叉验证折3 \\
2015--2020 & 2021 & 交叉验证折4 \\
2015--2021 & 2022 & 交叉验证折5（模型选择） \\
\hline
2015--2021 & 2023 & 测试（时间外评估） \\
2015--2021 & 2024 & 测试（时间外评估） \\
\hline
\end{tabular}
\end{table}

最终在2023和2024两个测试集上分别报告性能，再汇总2023--2024结果。所有数据预处理（填充器、编码器、标准化器）仅在训练集上拟合。

\subsubsection{辅助实验：按建筑ID分组}

增加按OSEBuildingID进行的GroupShuffleSplit实验，训练集与测试集中的建筑完全不重叠，用于检验模型对``从未见过的新建筑''的泛化能力。

\subsection{模型选择与评估方法}

\subsubsection{模型序列}

模型按复杂度递增顺序评估：

\begin{enumerate}
    \item \textbf{基线A（全体均值）：}以训练集SiteEUI均值作为所有建筑的预测值；
    \item \textbf{基线B（建筑类型中位数）：}按建筑类型分组，以训练集各组SiteEUI中位数作为同类建筑的预测值；
    \item \textbf{ElasticNet：}带L1/L2正则化的线性模型，使用OneHotEncoder编码类别特征，随机搜索调参；
    \item \textbf{随机森林：}Bagging集成，缺失填补后OneHotEncoder编码，随机搜索调参；
    \item \textbf{XGBoost：}梯度提升树，OneHotEncoder编码，随机搜索调参；
    \item \textbf{LightGBM：}梯度提升树，支持原生类别特征处理，随机搜索调参。
\end{enumerate}

所有模型训练时对目标取$\log(1+\text{SiteEUI})$以缓解右偏分布的影响，预测后通过$\exp(\hat{y})-1$还原为原始尺度。

\subsubsection{评价指标}

在验证集与测试集上分别计算以下指标（均在原始尺度上报告）：
\begin{itemize}
    \item MAE（平均绝对误差）；
    \item RMSE（均方根误差）；
    \item $R^2$（决定系数）；
    \item 可选RMSLE（均方根对数误差）。
\end{itemize}

不将$R^2>0.75$设为刚性成功标准。更合理的目标是：最佳模型在测试集上的MAE应明显优于全体均值基线和建筑类型中位数基线（例如降低10\%--20\%）。

\subsubsection{分组报告}

评估时应分别报告以下子集的性能：
\begin{itemize}
    \item 2023年测试集指标；
    \item 2024年测试集指标；
    \item 不同建筑类型（办公楼、酒店、医院等）的指标；
    \item 高能耗建筑子集（SiteEUI高于同类型P75）的指标；
    \item 已见建筑（训练集中出现过的OSEBuildingID）与未见建筑的指标。
\end{itemize}

\subsection{SHAP可解释性分析}

仅对最佳模型和测试集进行SHAP分析：
\begin{enumerate}
    \item 计算测试集上的SHAP值；
    \item 输出SHAP全局重要性条形图；
    \item 绘制SHAP summary beeswarm图；
    \item 绘制BuildingAge与LogGFA的SHAP依赖图；
    \item 绘制典型高能耗建筑与低能耗建筑的对比瀑布图；
    \item 对高度相关特征进行分组分析，避免重复计数；
    \item 报告中统一使用``重要预测特征''表述，明确标注SHAP反映的是模型内部的关联模式，不等于因果关系，不使用``驱动因素''或``因果因素''等措辞。
\end{enumerate}

\subsection{异常高耗能建筑识别}

不再使用``残差$>0$''直接判定异常，改为三层递进筛选：

\subsubsection{第一层：预测残差异常}

计算各建筑的预测残差：
\[\text{Residual} = \text{ActualEUI} - \text{PredictedEUI}\]
再按建筑类型和年度进行标准化：
\[\text{StandardizedResidual} = \frac{\text{Residual} - \mu_{\text{type, year}}}{\sigma_{\text{type, year}}}\]
异常条件：标准化残差超过合理的统计阈值（如1.5倍标准差），或残差位于同类型同年度前10\%。

\subsubsection{第二层：同类基准超标}

计算相对于同类P75基准的超耗比例：
\[\text{BenchmarkExcessRatio} = \frac{\text{ActualEUI} - \text{BenchmarkP75}}{\text{BenchmarkP75}}\]
要求同时满足：超出同类型P75，且超耗比例达到合理分位数阈值。

\subsubsection{第三层：持续性验证}

检查2022--2024年中的持续表现：
\begin{itemize}
    \item 进入同类高耗能前25\%的年数；
    \item 残差进入同类前10\%的年数。
\end{itemize}
最终候选需同时满足三层条件：模型预测异常 + 同类基准超标 + 最近三年持续出现。

\subsection{建筑级记录聚合}

原始数据为``建筑—年份''粒度，最终需要每栋建筑一条优先级记录。按OSEBuildingID汇总：
\begin{itemize}
    \item \textbf{当前状态（2024年）：}ActualEUI、PredictedEUI、StandardizedResidual、BenchmarkExcessRatio、GHGIntensity；
    \item \textbf{持续性表现（2022--2024年）：}最近三年高耗能次数、最近三年平均残差；
    \item \textbf{建筑属性：}建筑名称、建筑类型、BuildingAge、GFA；
    \item \textbf{数据质量：}最近三年数据完整率。
\end{itemize}

每栋建筑最终仅保留一行。

\subsection{优先级评分体系}

\subsubsection{评分维度设计}

\begin{table}[H]
\centering
\caption{节能改造优先级评分维度}
\label{tab:scoring}
\begin{tabular}{p{2.5cm}p{1.2cm}p{2cm}p{6cm}}
\hline
\textbf{维度} & \textbf{权重} & \textbf{方向} & \textbf{计算方式} \\
\hline
标准化预测残差 & 25\% & 正向 & 同建筑类型、同年度的标准化残差，归一化到0--1 \\
\hline
同类基准超耗程度 & 20\% & 正向 & 超出同类型P75基准的比例（BenchmarkExcessRatio），归一化 \\
\hline
持续高耗能频率 & 15\% & 正向 & 2022--2024年中进入同类高耗能前25\%的年数比例 \\
\hline
理论超额能耗量 & 15\% & 正向 & $\max(\text{超额EUI}, 0) \times \text{建筑面积}$，归一化（不称为``预计节能潜力''，因未考虑改造成本、技术方案、实际节能率和投资回收期） \\
\hline
GHG排放强度 & 15\% & 正向 & 单位面积碳排放（GHGEmissionsIntensity），归一化 \\
\hline
建筑年龄 & 10\% & 正向 & BuildingAge标准化到0--1 \\
\hline
\end{tabular}
\end{table}

\subsubsection{数据质量折扣}

定义数据质量得分为：
\[\text{DataQualityScore} = 1 - \text{MissingRate}\]
最终得分：
\[\text{FinalScore} = \text{RawPriorityScore} \times \text{DataQualityScore}\]
数据缺失严重的建筑不会因某个极端值被错误排到前面。

\subsubsection{敏感性分析}

设置四组权重方案进行对比：

\begin{table}[H]
\centering
\caption{权重敏感性分析方案}
\label{tab:sensitivity}
\begin{tabular}{p{2.5cm}p{1.8cm}p{1.8cm}p{1.8cm}p{1.8cm}p{1.8cm}p{1.8cm}}
\hline
\textbf{方案} & 标准化残差 & 基准超耗 & 持续频率 & 超额能耗 & GHG强度 & 建筑年龄 \\
\hline
标准方案 & 25\% & 20\% & 15\% & 15\% & 15\% & 10\% \\
节能优先 & 35\% & 25\% & 15\% & 15\% & 5\% & 5\% \\
减碳优先 & 15\% & 10\% & 10\% & 15\% & 40\% & 10\% \\
大型建筑优先 & 20\% & 15\% & 10\% & 30\% & 15\% & 10\% \\
\hline
\end{tabular}
\end{table}

\subsubsection{排序稳定性验证}

由于没有真实改造结果标签，通过以下方式间接验证排序的稳健性：
\begin{itemize}
    \item 不同权重方案下Top-50候选清单的重合率与Jaccard相似度；
    \item Spearman秩相关系数与Kendall秩相关系数；
    \item Bootstrap重采样后每栋建筑进入Top-50的频率；
    \item 各建筑类型在Top-50中的占比是否合理（避免某类型垄断）；
    \item 与官方benchmark高耗能建筑列表的一致性检验。
\end{itemize}

\subsection{辅助聚类画像}

聚类不再负责判断``谁异常''（该任务已由三层异常识别完成），而是对高优先级候选建筑进行画像，辅助理解候选建筑的类型结构：

\begin{itemize}
    \item \textbf{聚类对象：}优先级排名较高的候选建筑（非全部建筑）；
    \item \textbf{聚类特征（仅连续变量）：}LogGFA、BuildingAge、BenchmarkExcessRatio、持续高耗能频率、GHGIntensity、理论超额能耗量；
    \item \textbf{建筑类型角色：}不参与K-Means距离计算，仅在聚类后用于解释各聚类的类型构成；
    \item \textbf{K值选择：}综合肘部法则、轮廓系数、Calinski-Harabasz指数确定；
    \item \textbf{画像命名示例：}``大型老旧高超耗办公楼''、``高排放高强度酒店''、``面积较小但连续异常的商业建筑''、``大体量高绝对节能规模建筑''。
\end{itemize}


\section{研究计划与可交付成果}

\subsection{总体技术路线}

本项目按照一条清晰主线组织工作，各阶段依次推进、前序输出为后序输入：

\begin{center}
\begin{tabular}{c}
\hline
\textbf{阶段一：数据审计与预处理} \\
$\downarrow$ \\
\textbf{阶段二：探索性数据分析} \\
$\downarrow$ \\
\textbf{阶段三：能耗基准回归预测} \\
$\downarrow$ \\
\textbf{阶段四：模型解释} \\
$\downarrow$ \\
\textbf{阶段五：异常高耗能识别与建筑级汇总} \\
$\downarrow$ \\
\textbf{阶段六：改造优先级排序与稳定性检验} \\
$\downarrow$ \\
\textbf{阶段七：辅助聚类画像} \\
$\downarrow$ \\
\textbf{阶段八：可视化与结果分析} \\
\hline
\end{tabular}
\end{center}

\subsection{各阶段任务说明}

\subsubsection{阶段一：数据审计与预处理}

\begin{itemize}
    \item \textbf{数据挖掘问题类型：}数据清洗、缺失值处理、异常值检测
    \item \textbf{输入：}2015--2024年原始CSV文件（建筑—年份粒度，约34,699条记录，47个字段）
    \item \textbf{核心处理：}
    \begin{enumerate}
        \item 跨年度字段审计：检查各年份字段名称一致性、单位一致性、建筑类型命名规范性和建筑ID稳定性
        \item 数据质量评估：统计各字段缺失率、重复记录比例、数值字段解析异常
        \item 目标泄漏字段识别与剔除：排除与SiteEUI存在确定性数量关系的同期派生字段（SourceEUI、SiteEnergyUse等）
        \item 目标变量清洗：删除SiteEUI缺失、非正值和录入错误记录，保留合理的高能耗极端值
        \item 缺失值处理：数值字段优先按同建筑类型中位数填补，仍缺失则用训练集整体中位数；类别字段填充为Unknown；同步生成缺失指示变量
        \item 异常值处理：YearBuilt大于DataYear的记录标记为异常，按建筑类型中位年份填补，保证建筑年龄非负
        \item 数据集划分：按年份划分训练集（2015--2021）、验证集（2022）和测试集（2023、2024）
    \end{enumerate}
    \item \textbf{输出：}清洗后的特征矩阵与目标向量、数据质量审计报告（字段映射表、各年字段覆盖率、缺失值统计、重复记录报告）
\end{itemize}

\subsubsection{阶段二：探索性数据分析}

\begin{itemize}
    \item \textbf{数据挖掘问题类型：}描述性统计与可视化
    \item \textbf{输入：}阶段一清洗后的数据
    \item \textbf{核心处理：}
    \begin{enumerate}
        \item SiteEUI分布分析（原始尺度与对数变换后对比）
        \item 各建筑类型样本量与能耗水平对比（箱线图）
        \item 数值特征相关性分析（热力图）
        \item 建造年份与能耗关系分析（散点图与趋势线）
        \item 缺失值分布可视化
        \item 基准数据集匹配与超耗比例分布
    \end{enumerate}
    \item \textbf{输出：}EDA可视化图表集
\end{itemize}

\subsubsection{阶段三：能耗基准回归预测}

\begin{itemize}
    \item \textbf{数据挖掘问题类型：}监督学习—回归
    \item \textbf{输入：}训练集（2015--2021）、验证集（2022）、测试集（2023、2024）
    \item \textbf{核心处理：}
    \begin{enumerate}
        \item 基线构建：全体均值基线（Baseline A）和建筑类型中位数基线（Baseline B）
        \item 特征工程：构造BuildingAge、LogGFA、面积比率、混合用途标识、缺失指示变量等建筑属性特征；类别特征采用OneHotEncoder编码（禁止对建筑类型使用LabelEncoder）
        \item 模型训练与调优：按复杂度递增顺序评估ElasticNet、随机森林、XGBoost、LightGBM四种模型，采用按年份扩展窗口滚动交叉验证（2015--2017$\rightarrow$2018，\ldots，2015--2021$\rightarrow$2022），以负MAE为目标进行随机搜索调参；训练时对目标取对数缓解右偏，评估时还原为原始尺度
        \item 泛化能力检验：主实验采用时间外测试（2023、2024年）；辅助实验按建筑ID分组交叉验证以检验对未见建筑的泛化能力
        \item 设置两组对比实验：模型A不使用ENERGY STAR Score（主结果），模型B使用ENERGY STAR Score（扩展对比）
    \end{enumerate}
    \item \textbf{输出：}最优回归模型，各模型在验证集和测试集上的MAE、RMSE、$R^2$指标（按年份、建筑类型、高能耗子集、已见/未见建筑分组报告）
\end{itemize}

\subsubsection{阶段四：模型解释}

\begin{itemize}
    \item \textbf{数据挖掘问题类型：}模型可解释性分析
    \item \textbf{输入：}最佳模型、测试集数据
    \item \textbf{核心处理：}
    \begin{enumerate}
        \item 计算测试集SHAP值
        \item 生成全局特征重要性排名（条形图、beeswarm图）
        \item 绘制关键特征（BuildingAge、LogGFA）的SHAP依赖图
        \item 绘制典型高能耗与低能耗建筑的对比瀑布图
        \item 对高度相关特征进行分组解释
    \end{enumerate}
    \item \textbf{输出：}SHAP可解释性图表集；报告中统一使用``重要预测特征''表述，明确标注SHAP反映模型关联而非因果关系
\end{itemize}

\subsubsection{阶段五：异常高耗能识别与建筑级汇总}

\begin{itemize}
    \item \textbf{数据挖掘问题类型：}异常检测（基于残差分析和统计基准）
    \item \textbf{输入：}模型预测结果、原始SiteEUI、建筑类型—年份基准数据
    \item \textbf{核心处理：}
    \begin{enumerate}
        \item 第一层——预测残差异常：按建筑类型和年度计算标准化残差，标记标准化残差显著偏高或位于同类型同年度前10\%的建筑—年份记录
        \item 第二层——同类基准超标：计算实际SiteEUI超出同类型P75基准的比例，标记超耗比例显著的记录
        \item 第三层——持续性验证：检查2022--2024年中各建筑进入高耗能前25\%和残差前10\%的年数
        \item 建筑级聚合：按OSEBuildingID将多年记录汇总为建筑级结果（2024年当前状态 + 2022--2024年持续性表现），每栋建筑仅保留一行
    \end{enumerate}
    \item \textbf{输出：}异常高耗能建筑候选列表、建筑级汇总表（含当前状态和持续性指标）
\end{itemize}

\subsubsection{阶段六：改造优先级排序与稳定性检验}

\begin{itemize}
    \item \textbf{数据挖掘问题类型：}多指标综合排序评估
    \item \textbf{输入：}建筑级汇总表
    \item \textbf{核心处理：}
    \begin{enumerate}
        \item 6维度加权评分：标准化预测残差、同类基准超耗程度、持续高耗能频率、理论超额能耗量、GHG排放强度、建筑年龄
        \item 数据质量折扣：最终得分乘以数据完整率，防止缺失严重的建筑因极端值错误排前
        \item 多权重方案对比（标准/节能优先/减碳优先/大型建筑优先）
        \item 排序稳定性验证：不同方案下Top-N重合率与Jaccard相似度、Spearman与Kendall秩相关系数、Bootstrap重采样后各建筑入选频率、各建筑类型候选占比合理性
    \end{enumerate}
    \item \textbf{输出：}全量建筑优先级排名表、Top-N候选改造清单、权重敏感性分析报告、Bootstrap入选频率统计
\end{itemize}

\subsubsection{阶段七：辅助聚类画像}

\begin{itemize}
    \item \textbf{数据挖掘问题类型：}聚类分析（辅助性）
    \item \textbf{输入：}高优先级候选建筑的特征数据
    \item \textbf{核心处理：}
    \begin{enumerate}
        \item 聚类对象限定为优先级排名较高的候选建筑（非全部建筑），聚类不再负责判断``谁异常''
        \item 聚类特征仅使用连续数值变量（建筑面积对数、建筑年龄、基准超耗比例、持续高耗能频率、GHG强度、理论超额能耗量），建筑类型不参与距离计算，仅用于聚类结果的事后解释
        \item 综合肘部法则、轮廓系数、Calinski-Harabasz指数确定K值
        \item 为各聚类计算特征均值画像并命名（如``大型老旧高超耗办公楼''、``高排放高强度酒店''等）
    \end{enumerate}
    \item \textbf{输出：}聚类画像报告与PCA降维可视化图表
\end{itemize}

\subsubsection{阶段八：可视化与结果分析}

\begin{itemize}
    \item \textbf{数据挖掘问题类型：}结果呈现与综合分析
    \item \textbf{输入：}前序各阶段输出
    \item \textbf{核心处理：}汇总各阶段关键结果，生成综合分析可视化，撰写分析报告
    \item \textbf{输出：}Jupyter Notebook交互式分析报告、项目结题报告
\end{itemize}

\subsection{可交付成果}

\begin{enumerate}
    \item 完整的Python源码（模块化结构，各阶段独立可运行，通过主控脚本串联为完整流程）
    \item Jupyter Notebook交互式分析报告
    \item 训练好的模型文件
    \item 数据质量审计报告（含字段映射、各年覆盖率、缺失值统计、重复记录核查）
    \item 可视化图表集合（EDA图表、SHAP可解释性图表、聚类画像图表）
    \item 优先级排序输出（全量排名表、Top-N候选清单、敏感性分析报告）
    \item 开题报告与项目结题报告
\end{enumerate}



\section{项目特色与预期改进}

\subsection{项目特色}

\begin{enumerate}
    \item \textbf{清晰的数据到决策分析主链：}本课题不堆砌方法，而是以一条清晰主线组织工作——``跨年度数据审计 $\rightarrow$ 能耗基准预测 $\rightarrow$ 异常高耗能识别 $\rightarrow$ 建筑级汇总 $\rightarrow$ 改造优先级排序 $\rightarrow$ 辅助聚类画像''。每个环节有明确的输入输出和质量检查节点，聚类降级为辅助画像而非核心目标。区别于大多数仅训练单一模型或简单排列方法的课题。

    \item \textbf{按年份滚动验证与双重泛化检验：}主实验采用按完整年份的扩展窗口滚动验证（2015--2017$\rightarrow$2018至2015--2021$\rightarrow$2022），更真实地模拟``用历史数据估计基准''的场景；辅助实验采用按OSEBuildingID的GroupShuffleSplit，检验模型对从未见过的新建筑的泛化能力。两种验证策略互补，全面评估模型性能。

    \item \textbf{三层递进式异常识别机制：}不再使用``残差$>0$''直接判定异常，而是采用三层递进筛选——（1）模型预测残差异常（同类型同年度标准化残差$>1.5$或前10\%）；（2）同类P75基准超标；（3）2022--2024年持续性验证。三层交叉验证可显著降低单一方法的误判率。

    \item \textbf{建筑级记录聚合与持续性评估：}原始数据为``建筑—年份''粒度，本研究按OSEBuildingID将多年记录聚合为建筑级结果（2024年当前状态 + 2022--2024年持续性表现），每栋建筑仅保留一行。持续高耗能频率和最近三年平均残差等指标使优先级排序更加稳健。

    \item \textbf{可配置的优先级评分与全面的稳定性检验：}设计6维度加权评分框架（标准化残差25\%、基准超耗20\%、持续频率15\%、超额能耗量15\%、GHG强度15\%、建筑年龄10\%），引入数据质量折扣因子防止缺失严重的建筑因极端值错误排前；通过Bootstrap、Jaccard相似度、Spearman/Kendall秩相关系数等多维度检验排序稳定性。支持4套权重方案对比（标准/节能优先/减碳优先/大型建筑优先）。

    \item \textbf{全面的评估基线与分组报告：}以全体均值、建筑类型中位数作为基线，当且仅当复杂模型显著优于简单基线时才采信其结果；评估时按年份、建筑类型、高能耗子集、已见/未见建筑分组报告性能指标。

    \item \textbf{高度模块化的代码结构：}7个独立Python脚本，每个模块职责单一、输出明确，可独立运行或通过主控脚本一键执行。所有预处理拟合严格限定在训练集，避免数据泄漏。

    \item \textbf{完善的跨年度数据审计：}在正式建模前对各年份数据进行系统性审计（字段名称一致性、单位一致性、建筑类型命名一致性、ID稳定性、重复记录、数值解析、字段覆盖率），生成字段映射表和覆盖率统计等审计文件，确保数据质量可追溯。
\end{enumerate}

\subsection{与现有工作的区别}

\begin{enumerate}
    \item 现有建筑能耗benchmarking项目多使用ENERGYSTAR Score作为主要依赖（该分数本身由当年能耗数据计算），本研究设置两组实验——主模型（模型A）不使用ENERGYSTAR Score，仅保留ES\_Missing指示变量，模型B作为扩展对比，更真实地模拟``未获评建筑''的预测场景并防止模型性能因评级信息而虚高；

    \item 将基准性能范围数据（Benchmarking Performance Ranges by Building Type）作为异常识别第二层（同类基准超标）的参照系，而非仅用于EDA展示；

    \item 在优先级排序中采用三层递进筛选（预测残差异常 + 同类基准超标 + 持续性验证），而非简单的正残差截断；同时引入数据质量折扣因子，避免因单一极端值导致的错误推荐；

    \item 不同建筑类型天然能耗水平差异巨大（如医院vs仓库），直接用绝对SiteEUI排序会导致某些类型垄断Top-50，通过同类型超耗比例替代绝对能耗，使排序更加公平；

    \item 将聚类从``识别异常''的核心角色降级为``候选建筑画像''的辅助角色，聚类对象仅为高优先级候选建筑，使用连续数值特征计算距离，建筑类型仅用于事后解释，逻辑更加清晰。
\end{enumerate}

\subsection{待改进方面}

\begin{enumerate}
    \item \textbf{数据扩充可能性：}西雅图公开数据未包含墙体保温、窗户类型等工程细节，若引入EnergyPlus模拟或结合气候数据（如采暖/制冷度日数），可显著提升预测精度。

    \item \textbf{评分权重可进一步优化：}当前权重为手动设定并通过敏感性分析验证，如引入实际改造成本数据（如ECM成本面板），可将主观排序目标转化为成本效益优化问题。

    \item \textbf{因果关系需慎重声明：}SHAP显示的特征重要性属于模型内部关联，不代表真实因果关系；在实际政策决策中需结合领域知识谨慎解读。报告中统一使用``重要预测特征''表述。

    \item \textbf{聚类稳定性：}K-Means对初始中心敏感且假设各向同性，如使用DBSCAN或谱聚类可能发现非凸形状的聚类结构。但由于聚类已降级为辅助画像，该问题对核心结论的影响有限。

    \item \textbf{滞后特征未纳入：}当前主模型仅使用静态建筑属性和年度背景变量，未加入上一年SiteEUI、近三年均值、同比变化率等滞后特征。若后续研究需要严格的时间序列预测，可在此基础上扩展。
\end{enumerate}


% ==================== 参考文献 ====================
\begin{thebibliography}{99}

\bibitem{unep2020global}
United Nations Environment Programme,
``2020 Global Status Report for Buildings and Construction: Towards a Zero-emission, Efficient and Resilient Buildings and Construction Sector,''
UNEP, Nairobi, 2020.

\bibitem{seattle_benchmarking}
City of Seattle,
``Building Energy Benchmarking Data, 2015--Present,''
Office of Planning \& Community Development,
\url{https://data.seattle.gov/dataset/Building-Energy-Benchmarking-Data-2015-Present/bfsh-8r6z}.
[Accessed: 2026-07].

\bibitem{ashrae_energy}
ASHRAE,
``ASHRAE Guideline 14-2014: Measurement of Energy, Demand, and Water Savings,''
American Society of Heating, Refrigerating and Air-Conditioning Engineers, Atlanta, GA, 2014.

\bibitem{rf_breiman}
L.~Breiman,
``Random Forests,''
\textit{Machine Learning}, vol.~45, no.~1, pp.~5--32, 2001.

\bibitem{xgboost_chen}
T.~Chen and C.~Guestrin,
``XGBoost: A Scalable Tree Boosting System,''
in \textit{Proceedings of the 22nd ACM SIGKDD International Conference on Knowledge Discovery and Data Mining},
San Francisco, CA, USA, 2016, pp.~785--794.

\bibitem{lgbm_ke}
G.~Ke, Q.~Meng, T.~Finley, \textit{et al.},
``LightGBM: A Highly Efficient Gradient Boosting Decision Tree,''
in \textit{Advances in Neural Information Processing Systems 30 (NIPS 2017)},
Long Beach, CA, USA, 2017, pp.~3146--3154.

\bibitem{shap_lundberg}
S.~M.~Lundberg and S.-I.~Lee,
``A Unified Approach to Interpreting Model Predictions,''
in \textit{Advances in Neural Information Processing Systems 30 (NIPS 2017)},
Long Beach, CA, USA, 2017, pp.~4765--4774.

\bibitem{ahp_saaty}
T.~L.~Saaty,
``How to make a decision: The Analytic Hierarchy Process,''
\textit{European Journal of Operational Research}, vol.~48, no.~1, pp.~9--26, 1990.

\bibitem{energystar}
U.S.~Environmental Protection Agency,
``ENERGY STAR Portfolio Manager: Technical Reference,''
EPA, 2018.
\url{https://www.energystar.gov/buildings/portfolio-manager}. [Accessed: 2026-07].

\end{thebibliography}

\end{document}